\documentclass[11pt]{article}
\usepackage[T1]{fontenc}
\usepackage[utf8]{inputenc}
\usepackage{lmodern,amsmath,amssymb,amsthm}
\usepackage[margin=1in]{geometry}
\usepackage[hidelinks]{hyperref}
\newtheorem{theorem}{Theorem}
\newtheorem{lemma}[theorem]{Lemma}
\newcommand{\Z}{\mathbb Z}
\newcommand{\N}{\mathbb N}
\newcommand{\C}{\mathbb C}
\newcommand{\R}{\mathbb R}
\newcommand{\E}{\mathbb E}
\newcommand{\GI}{\mathrm{GI}}
\title{Assembly of the precise inverse-theorem input}
\author{Author successor: /root/astra\_v14\_baseline\_theta}
\date{7 November 2026}
\begin{document}
\maketitle

\noindent\textbf{Status.} This successor assembles the exact local proofs
identified below and retains their separate independent decisions.
The preceding assemblies' local splitting, selection, analytic and
finite-family joins have conditional acceptance. The corrected base,
degree conversions, rank iteration and final Section 12 and Section 13
units have their own scoped decisions. The local converse and plain
additive-quadruple argument have also been independently accepted.
The actual Sections 9--11 reconstruction has conditional acceptance.
The one-variable Appendix D/F packet has bounded decisions. This
successor connects its affine-offset transfer and the multivariable
E.11/E.13 uses to the independently reviewed local multifactor
supplier. Its upstream and downstream decisions are separate.
The new composition and remaining source-coverage review must
still be assessed; no release approval is presumed here.
This document does not convert scoped decisions into unconditional GI
or whole-source acceptance. Exact versions, decisions and open
interfaces are bound in the accompanying records.

\section{Exact target and conventions}

The target, denoted $\GI(k)$, is the following unchanged linear formulation.
For each fixed integer $k\geq3$ and each $0<\delta\leq1$, there are a finite
family $\mathcal M_{k,\delta}$ of nilmanifolds $G/\Gamma$, each with connected
simply connected nilpotent group of step at most $k$ and a discrete cocompact
lattice, fixed smooth metrics, and constants $C,c>0$, such that every
$1$-bounded $f:[N]\to\C$ with $\|f\|_{U^{k+1}[N]}\geq\delta$ correlates with
a $1$-bounded test $F$ of Lipschitz constant at most $C$ along a linear orbit:
\[
 \left|\E_{n\in[N]}f(n)F(g^nx)\right|\geq c.
\]
The family and constants depend on $k,\delta$, not on $N,f$ or the orbit
parameters. Strict source thresholds can be used with $\delta/2$ to give
this weak-inequality version. No effective numerical bounds are asserted.

For $d\geq2$, the interval norm is the norm of the zero extension to a
cyclic group of sufficiently large order divided by the norm of the
interval indicator, as in the published GTZ definition on p.~1232.
We use the same normalization at every inverse step. On a finite abelian
group $A$, write
\[
 \|v\|_{U^d(A)}^{2^d}
 =\E_{x,t_1,\ldots,t_d\in A}
   \prod_{\omega\in\{0,1\}^d}\mathcal C^{|\omega|}
     v(x+\omega\cdot t),
\]
where $\mathcal C$ is complex conjugation. Our correlation convention is
$|\E f\psi|$; conjugating the test function converts the other common
convention without changing its degree or bounds.

All limit nilcharacters in the assembly are internal in the explicitly
specified sense of the local Lemma 13.2 and tautification units: the
vertical identity holds for all limit group elements and all limit points,
with the natural extension of a standard vertical character. Every
constructor used below must supply that domain, rather than only
standard-point equivariance. This is an explicit application check, not a
consequence of bounded Lipschitz continuity alone.

The manuscript currently requires $\GI(20)$ and $\GI(39)$. The proposed
relative inverse replacement also requires $\GI(r-1)$ for each fixed
$r=\max(s,4)$ chosen before the norm threshold and asymptotic variable.
A polynomial orbit without a same-step linearization or an
$N$-dependent family does not satisfy this target.

\section{Intermediate interfaces that must be supplied}

The following lists the available components and the remaining
mathematical premises of the assembly argument. A local decision applies
only to its stated scope; the incoming source contracts remain explicit.

\begin{enumerate}
\item The compatible bounded polynomial inverse theorem of order $2$,
proved in \texttt{gi-base-002} and independently accepted relative to
its retained published forward bracket theorem and Bohr-set inputs.
It gives the fixed finite family and normalized interval $U^3$ norm.
The needed trivial-subgroup specialization of the false quadratic
extension is proved locally, and the periodic connected polynomial
construction replaces the affected skew-iterate route.
\item The nilcharacter approximation/selection interface: from positive
correlation with a bounded degree-$d$ nilsequence, obtain a degree-$d$
unit-vector nilcharacter and a scalar component giving positive correlation.
For a dense family of parameters, make the choices in a limit fashion and,
after passage to a dense subfamily if necessary, keep the finite dimensions,
complexities, bounds and positive correlation standard and uniform.
Section 4 below gives the finite-net and unit-vector construction used
at this interface, with conditional independent acceptance against the
previous finite inverse family. Later selection applications retain
their own source contracts.
\item Linearization at inverse order $s$: for a dense family of degree
$s-1$ nilcharacters $\chi_h$ correlated with $\Delta_hf$, obtain one
multidegree $(1,s-1)$ nilcharacter $\chi(h,n)$ with which $\Delta_h f$
has degree-at-most-$s-2$ correlation for many $h$.
This is the precise Theorem 7.1 interface. The modified Proposition 7.3
is supplied by \texttt{local-proof-sol-lp02-03-r3}; the quotient and
rank-lower comparison by \texttt{lp04-001}; the final orbit by the
accepted Section 12 successor; and the degree conversions, unit-vector
recovery, finite rank iteration and rank-zero endpoint by
\texttt{rank-iteration-001}. The frozen Sections 9--11 reconstruction
supplies an independently reviewed proof of the one-step frequency premise at
every newly constructed rank family, conditional on its precise
one-variable D and F contracts. Its theorem takes the same uniform
unit-vector degree--rank family and accumulated multidegree factor as
the rank iteration. It reconstructs all core/linear/independent data
anew; none is inherited. Its fixed unit padding is removed by projecting
the first factor of the entire LP04 matrix, so the matrix is exactly
$\chi_h\otimes\overline{\zeta_h}$ required by the next rank step.
The Section 12 successor supplies the same fixed coordinate isometry
as the earlier version named in the frozen rank packet; its only
declared change makes the zero-index groups explicit. The new binding
is recorded here without rewriting that earlier freeze. All these
steps retain the full-limit vertical identities and uniform bounds.


For clarity, the invariant concerns represented families: ordinary
unit-vector tests with exact global vertical identities, fixed standard
bounds and finitely many standard models and characters. The initial
positive-degree selection is made among these representations, by the
parameter-level selection argument for the preceding inverse order;
it does not upgrade an arbitrary weak nilcharacter or preserve a
prescribed symbol. Degree-to-rank conversion retains the represented
function and restricts its character, so this invariant starts at
$(m,m)$, where $m=s-1$, with accumulated factor $A=1$.

Distinguish the fixed inclusion $J$ that pads the incoming $\chi_h$
in Sections 9--11 from the final Section 12 inclusion $I$, for which
$Z(h,n)=I\zeta_h(n)$ on the retained shift set. The full ordinary-top
character in Section 12 is chosen before the LP04 comparison and
restricts to its cancelling rank-top character. The comparison thus
supplies the entire vector
\[
 W_h=(J\chi_h)\otimes\overline{\zeta_h}
\]
on one common model of degree-rank at most $(m,R-1)$. If $QJ=\mathrm{id}$
is the coordinate projection on the first factor, then
\[
 (Q\otimes\mathrm{id})W_h
   =\chi_h\otimes\overline{\zeta_h}=:V_h.
\]
This fixed bounded linear map preserves the common model, internal
parameter dependence and standard Lipschitz bound. The projection
recovering $\zeta_h$ from $Z$ is instead the one associated with $I$.
In particular, this conclusion concerns every matrix entry, not a
separately selected scalar entry.

The identity
$\chi_{h,a}=\sum_b\zeta_{h,b}(\chi_{h,a}\overline{\zeta_{h,b}})$
uses the unit norm of the whole vector $\zeta_h$ and therefore uses
exactly $V_h$ in the correlation recovery. The finite ordinary restart
dictionary preserves the global unit and vertical identities after
internal selection; the ordinary tests chosen in the proof of E.5 are
extended only after their finite selection. Tensoring the accumulated
factor with $Z$ preserves its unit norm and multidegree $(1,m)$.
Sections 9--11 then reconstruct all frequency data for this new family.
No ordinary degree drop is claimed before the rank endpoint: at rank
zero the degree-rank axiom makes the remaining factor ordinary degree
at most $m-1$, which is then absorbed into the lower-degree test.
The geometric, analytic and source premises of all these component
constructions remain those of their stated contracts.

The centered fractional-part convention in this chain is
$(-1/2,1/2]$, a fixed fundamental cube permitted by Section 12.
The reconstruction's acceptance remains conditional on its suppliers;
the affine-offset coefficient transfer is supplied by the local multifactor lemma, as connected below.
\item Symmetrization at inverse order $s$: from that input obtain one
degree-$s$ nilcharacter $\Theta$ and a bounded nilsequence $\Psi(h,n)$
supported on
\begin{equation}\label{eq:J}
 J=\{(i,j):i+j\leq s-1\}\ \cup\ \{(i,s-i):2\leq i\leq s\},
\end{equation}
such that the relevant $\chi(h,n)$ is a bounded finite linear combination
of components of
$\Theta(n+h)\otimes\overline{\Theta(n)}\otimes\Psi(h,n)$.
This is the required Theorem 7.4 output, supplied by the locally authored
Section 13 proof with its scoped independent decision. It takes the
degree-at-most-$s-2$ correlation actually supplied above; the printed
strict hypothesis is not used. Its proof binds the corrected
multilinearization, literal block symmetry, local Lemma 13.2 and later
shift/sign calculations. It also retains the precise stronger
Proposition 8.4, E.11 and E.13 applications. The precise plain
Proposition 8.4 input is supplied by the accepted local
\texttt{plain-quadruples-001}: select the same scalar coordinate of
$\chi(h,\cdot)\otimes\phi_h$ and retain that entire lower factor.
For $H\subset[-N,N]$ of density $\alpha$, correlation $\delta$ and
global scalar bound $B$, that proof gives standard positive density
and correlation depending only on $\alpha,\delta,B$. Its output has
signs $++--$ and shifts in the second and fourth factors by
$h_1-h_4$, exactly as used in Section 13. The functions have global
polynomial representations; the interval cutoff belongs to $f$.
The proof removes its remaining bounded cutoff by an inequality,
not by an asserted identity. Finite component selection and the
internal version preserve standard uniform constants.
This is an explicit substitution of the supplying proof for the
published Proposition 8.4 invocation in the frozen Section 13 proof,
lines 88--125. That earlier file remains unchanged and does not itself
refer to the later local proof. Its intervals $I_1=[-N,N]$ and $J=[N]$,
scalar component, global test domain, and quadruple count on the
$N^3$ scale are exactly the local theorem's data. The reparametrization
$(h_1,h_2,h_3,h_4)=(h_0+a,h_0+b,h_0+a+b,h_0)$ is therefore unchanged;
the subsequent pigeonhole has $h_0\in I_1$ and $a,b\in[-2N,2N]$.
The E.11 and E.13 inputs require the
actual multivariable Appendix D.5 factorization and all-box reconstruction;
a one-variable D.4 proof does not supply that premise. This obligation
is supplied by the local multifactor theorem and its two independent decisions, as connected below.
\item Splitting for a union of two finite downsets: a fixed bounded
nilsequence on their union is uniformly approximable, to any fixed positive
error, by a finite sum of products of nilsequences supported on the two
downsets, with fixed finite complexity and bounds. This is the E.4
application used explicitly in the next section.
\item The converse correlation bound: a bounded scalar function with
positive correlation against a bounded degree-$d$ nilsequence has positive
$U^{d+1}[N]$ norm, with a bound for each fixed complexity and correlation.
The local \texttt{converse-001} proves the required implication using
the published compact cube constraint, an explicitly normalized smooth
average, a uniform Fourier bound and Gowers--Cauchy--Schwarz. It uses
accepted exact same-step linearization for polynomial tests. It does not
import the disputed averaging-domain transition in Green--Tao 11.2 or
invoke the inverse theorem being proved. Its exact source application
and independent decision are bound separately.
\item Countability, up to filtered isomorphism, of the fixed rational
filtered nilmanifolds of bounded degree, and equivalence of fixed smooth
metrics on a compact nilmanifold. The published countability assertion is
on p.~1250 with its Mal'cev-coordinate justification. Its exact scope must
include the taut rational objects constructed by this project.
\item The local increment linearization theorem: an exact polynomial
nilsequence of degree $s$ has a same-step linear representation on fixed
geometry with a fixed Lipschitz comparison. This is supplied by the
separately authored and independently accepted $P_1\rtimes\R$ unit,
rather than by the defective full-$P_0$ suspension in Appendix C.
\end{enumerate}

The increment unit supplies the exact sequence and finite-family join
used here. It does not assert the stronger clause in the original
Proposition C.2 about one universal projection into the original,
unnormalized nilmanifold for every polynomial: bounded normalization is
absorbed in the test function $F_q$. That distinction changes no part of
the target in section 1.

The successor connections to the frozen multifactor proof are as follows.
They select a supplying proof for the existing conclusion and preserve
the old component files. Their independent decisions are recorded
separately when available.

First, the one-variable weak factorization uses a progression
$A+b[M]$ with $|A|=O(N)$, standard bounded nonzero integer $b$, and
$M\asymp N$. Apply the local inverse-affine lemma with one variable
to a real lift of the horizontal polynomial. Its input is the positive
binomial coefficient bound on $m\mapsto f(A+bm)$. In the inverse
substitution, the coefficient from input degree $j$ to positive output
degree $i$ has size $O(N^{j-i})$ and a standard denominator independent
of $A$. One standard integer multiplier clears the integer portions;
the residual $O(N^{-j})$ becomes $O(N^{-i})$. Thus the same original
horizontal character, multiplied by a standard positive integer,
has the required ambient coefficient bound. It remains nonzero.
No bound on the constant coefficient is asserted or needed. This is
the exact proof-level transfer at DF002 lines 132--145, distinct from
its right-coset periodicity argument.

Second, at Section 13 lines 327--343 take the local multifactor bias
theorem with dimension $s+1$, degree $s$, and its scale equal to $8N$.
The input is the actual reflected unit nilcharacter
$\sigma_\varepsilon$ and reflected convex polytope obtained by the
orthant selection. Every positive progression is the full $[8N]$;
the auxiliary lower-degree test is constant. Its conclusion supplies
a globally represented lower-degree nilsequence agreeing with
$\sigma_\varepsilon$ on all of $[-8N,8N]^{s+1}$.
The inverse reflection is an integer linear substitution, so the
same full-box conclusion holds for $\sigma$. This does not extend
an identity from only the selected positive orthant by assumption.

Third, after the displayed diagonal pullback in Section 13, the
actual unit tensor-product representative $\xi(h,n)$ of the
antisymmetric symbol satisfies $s![\xi]=0$ on $[-8N,8N]^2$.
By the definition of the zero symbol, the coordinates of
$\xi^{\otimes s!}$ are lower-degree nilsequences on that box;
their finite product of models is one common lower-degree model.
Apply the local torsion theorem in dimension two, degree $s$, scale
$8N$, and positive standard tensor power $s!$. It gives a globally
represented degree-at-most-$(s-1)$ nilsequence equal to $\xi$ on
that same symmetric box. Only afterward restrict to
$[-N,N]\times[N]$, as in Section 13 lines 345--363.

There is also a distinct converse use at Section 13 lines 298--305.
For $s\geq4$, put $d=s-3\geq1$ in the accepted local converse and
apply it to the uniformly selected lower-degree test correlated
with the chosen coordinate of $\alpha'$. This gives the positive
$U^{s-2}[N]$ bound there, uniformly after the stated selections.
For $s=3$, a degree-zero nilsequence is constant: a fixed positive
correlation with a bounded constant directly implies a fixed
positive $|\mathbb E_{n\in[N]}\alpha'(n)|$: if the correlation is
at least $c>0$ and the constant has modulus at most $M$, this mean
has modulus at least $c/M$. Zero extension has group mean
$(N/|Q|)\mathbb E_{[N]}\alpha'$, and the interval indicator has
$U^1$ quantity $N/|Q|$, so normalization gives exactly that interval
mean. No degree-one test is inserted into this case.

Closure of fixed-degree nilsequences under conjugation, scalar
multiplication and finite products uses the actual product nilmanifold and
product filtration. The top degree is the maximum of the input degree
bounds. This elementary object identification is also required when it is
used in the induction below.

\section{Correct splitting of the symmetrization remainder}

We first give the lattice detail needed by the splitting application.
It uses the integral Hall-coordinate lemma in the independently reviewed
LP04 unit, not a general assertion that disjoint rational Lie subgroups
automatically separate points after quotienting by a lattice.

\begin{lemma}[Splitting on two finite multidegree downsets]
Let $J_1,J_2\subset\N^a$ be finite downsets. Every bounded Lipschitz
nilsequence of multidegree supported on $J_1\cup J_2$ is uniformly
approximable, to each fixed positive error, by a finite sum of products
of nilsequences supported on $J_1$ and $J_2$. The construction uses
fixed geometries and bounds before the polynomial coefficients or the
interval length vary, for each fixed input geometry and test bound.
\end{lemma}
\begin{proof}
Use bounded tautification to set the polynomial's value at zero to the
identity and restrict the ambient group to the group generated by its
positive coordinate groups. For each positive $j\in J_1\cup J_2$ choose
a finite generating set of $\Gamma\cap G_j$. Its real generated group
is $G_j$, by the rational subgroup assertion. Introduce one formal
letter of weight $j$ for each such generator. Take the free real Lie
algebra on these letters and discard every individual-letter multidegree
whose total vector weight is outside $J_1\cup J_2$. This is an ideal:
the complement of a downset is upward closed. Every letter has positive
total weight, so the surviving length is bounded by
$\max_{j\in J_1\cup J_2}|j|$.

Let $U$ be the exponential group and $\Lambda$ the group generated by
the exponential letters. The Hall-coordinate lemma proves that
$\Lambda$ is a lattice and that all weight filtration groups are
rational. Sending each letter to its chosen original lattice generator
gives a filtered Lie homomorphism $\phi:U\to G$ with
$\phi(\Lambda)\subseteq\Gamma$. It maps $U_j$ onto $G_j$: the letters
of weight exactly $j$ already generate $G_j$ as a real group. Lift each
Taylor coefficient of the normalized original polynomial through this
surjection. The multivariate Taylor theorem, B.9, makes their ordered
binomial product an adapted polynomial $\widetilde g$ with
$\phi\widetilde g=g$. Pulling the test back along
$U/\Lambda\to G/\Gamma$ therefore preserves the sequence exactly.
Its Lipschitz bound increases by a fixed geometric factor only.

For $\nu=1,2$ let $I_\nu$ be the connected ideal subgroup spanned by
the surviving Hall words whose weight is outside $J_\nu$. The quotient
$U/I_\nu$ has degree support in $J_\nu$. The integral Hall lemma says
that its quotient lattice is obtained by deleting precisely those
coordinates. No surviving Hall word is deleted by both quotients, since
its weight belongs to $J_1\cup J_2$. Consequently the map
\[
 U/\Lambda\longrightarrow
 U/(I_1\Lambda)\ \times\ U/(I_2\Lambda)
\]
is injective. In detail, if two cosets have equal images, their difference
has integral coordinates in each projected Hall system. The union of
the retained coordinate sets is the full Hall system, so every coordinate
of the difference is integral. The difference is in $\Lambda$.
This checks the lattice condition as well as the Lie algebra condition.

The pullbacks of the two algebras of Lipschitz functions on these compact
quotients contain the constants, are closed under conjugation, and
together separate points by that injection. The complex
Stone--Weierstrass theorem approximates the pulled-back test uniformly
by the algebra they generate. Each element of that algebra is a finite
sum of products with one factor from each quotient, since each of the
two individual algebras is closed under multiplication. Composing with
$\widetilde g$ gives the asserted splitting. For an internal bounded
Lipschitz test, first approximate by one standard Lipschitz test on the
fixed compact manifold, using Definition 5.1, and then perform this
standard finite approximation. For a fixed uniform Lipschitz and norm
bound, compactness of that class in the uniform norm supplies a finite
net and thus fixed finite bounds. Only this fixed-bound form is used
below. If the positive groups are trivial, the sequence is constant and
the assertion follows with one constant product.
\end{proof}

The coordinate order is $(h,n)$. Set
\[
 A=\{(0,j):0\leq j\leq s-1\},\qquad
 B=\{(i,j):j\leq s-2,\ i+j\leq s\}.
\]
Both are downsets and $A\cup B=J$. Indeed the only point of $J$ with
$j=s-1$ is $(0,s-1)$, covered by $A$. Every other point of $J$ has
$j\leq s-2$ and lies in $B$. Conversely a point of $B$ of total degree
$s$ has $i\geq2$, so belongs to the second part of \eqref{eq:J}; all its
other points lie in the first part. This proves equality in both directions.

Apply interface 5 to the actual $\Psi$ with these two downsets. A sequence
supported on $A$ is independent of $h$: its polynomial map has every
$h$-direction derivative in the trivial first-coordinate group.
It is a degree-at-most-$s-1$ function $\psi(n)$ of $n$ alone.
For a factor supported on $B$, fixing $h$ gives a polynomial nilsequence
in $n$ of degree at most $s-2$, since all higher $n$-direction groups are
trivial. If normalization of this section at fixed $h$ is needed, the
accepted bounded tautification theorem supplies it with uniform geometry
and parameter dependence.

Here is the needed approximation-to-correlation connection. After uniform
selection, suppose the preceding correlation is at least a standard
$\epsilon_0>0$ on a dense set of $h$, and the lower-degree test factors
have magnitude at most a standard $M$. Choose the uniform splitting error
less than $\epsilon_0/(2M)$, also absorbing the fixed finite component and
coefficient bounds into $M$. The error in each normalized correlation is
then at most $\epsilon_0/2$. A finite sum remains, so one summand has a
positive standard correlation on a dense subfamily by finite pigeonhole.
The same finite choice of components of $\Theta$ can be made on a dense
subfamily. The result is two fixed bounded scalar nilsequences
$\theta,\theta'$ of degree at most $s$ such that
\begin{equation}\label{eq:factorcorr}
 n\longmapsto f(n+h)\overline{f(n)}\theta(n+h)\theta'(n)
\end{equation}
is degree-at-most-$s-2$ biased for a dense set of $h$.
The $B$-factor is absorbed into that lower-degree test. The $A$-factor
$\psi(n)$ is multiplied into the unshifted component forming $\theta'$.
No scalar component is divided by, or assumed to have modulus one.

The printed source on p.~1272 uses the opposite coordinate restrictions
while claiming independence of $h$. The corrected downsets and the
approximation estimate above are the explicit replacement for that step;
its historical printed form is preserved in the findings record.

\section{Uniform nilcharacter selection and fixed geometries}

Here is the selection needed at the start of each induction step. Use
the finite-family inverse statement already obtained at the previous
order; the finite-family extraction below is performed after each order,
before proceeding to the next. This avoids assuming that an arbitrary
external assignment of nilmanifolds to parameters is an internal family.

On one fixed taut degree-$d$ nilmanifold, the top group is central and
its quotient by its lattice is a compact torus $T$. For each integral
character $\eta$ of $T$, a unit-vector smooth function with that vertical
character exists as follows. Cover the compact quotient by $T$ by finitely
many trivializing charts. Choose smooth real functions $\rho_j$ supported
inside the charts with $\sum_j\rho_j^2=1$, using a smooth partition of
unity and normalizing its vector by its positive Euclidean norm. In the
$j$th trivialization define the $j$th component by
$\rho_j(b)e(\eta(t))$, and extend it by zero off that chart. These
components are smooth, have the exact character, and their squared
magnitudes sum to one. The central rational quotient and its lattice
are justified by the accepted quotient geometry in the local Lemma 13.2
unit. This is the explicit construction used in published Lemma 6.4.

Components of such unit-vector functions, including frequency zero,
form a conjugation-stable multiplicative family: use tensor products
for multiplication and conjugation for conjugation. They contain
constants and separate points. Distinct points on the same $T$-fibre
are separated by an integral character, with a nonzero component of
the unit vector. Points on different fibres are separated by a smooth
frequency-zero scalar unit function pulled back from the base.
Stone--Weierstrass therefore gives uniform approximation of every
continuous test by a finite linear combination of these components.

For the fixed norm and Lipschitz bound in the previous inverse theorem,
Arzela--Ascoli gives a finite uniform net of tests. Apply the preceding
approximation to each element of that net with error smaller than a
fixed fraction of the positive inverse correlation. There are finitely
many models and net elements, so the total list of unit-vector test
maps, dimensions, frequencies, coefficient bounds and Lipschitz bounds
is fixed. Pigeonholing a scalar component of an approximating sum
preserves a fixed positive correlation. No component is assumed to
have modulus one; the entire vector has norm one.

For each finite interval and each of its good shifts, ordinary choice
selects one polynomial orbit and one item from this fixed finite list.
Taking the ultralimit of these finite assignments gives an internal
family. A finite pigeonhole step restricts to a dense set of shifts with
one common fixed model and fixed test map if required. The unit and
vertical identities hold at every representative and hence for all
limit arguments, with the natural extension of the chosen standard
character. This establishes the selection interface needed here from
the preceding finite inverse theorem. It does not replace separate
selection applications later in Sections 8--13, which remain with their
actual consumers.

For the countability interface, use the rational Mal'cev-coordinate
description of a lattice in a connected simply connected nilpotent Lie
group, as invoked in published Section 3 and p.~1250. In each finite
dimension the rational structure constants have countably many
possibilities. The rational subspaces defining the finitely many
filtration groups are specified by finite matrices with rational
entries, again a countable choice. In the rational structure determined
by the lattice, a finite list of rational logarithmic generators
specifies the lattice. Thus there are at most countably many such
filtered lattice data up to isomorphism. All local constructions here
have finite rational filtrations and discrete cocompact lattices, so
belong to this class. On each compact quotient any two fixed smooth
Riemannian metrics are bi-Lipschitz equivalent: compare their positive
quadratic forms on the compact unit tangent bundle. These are the exact
countability and metric comparisons used in the extraction below.

\section{The finite-group inequality needed for the last induction step}

\begin{lemma}\label{lem:CS}
For bounded complex functions $a,b$ on a finite abelian group $Q$ and an
integer $d\geq1$,
\[
 \E_{h\in Q}\|a(\cdot+h)b(\cdot)\|_{U^d(Q)}^{2^d}
 \leq
 \|a\|_{U^{d+1}(Q)}^{2^d}
 \|b\|_{U^{d+1}(Q)}^{2^d}.
\]
\end{lemma}
\begin{proof}
For $t=(t_1,\ldots,t_d)$ write
$C_a(x;t)=\prod_\omega\mathcal C^{|\omega|}a(x+\omega\cdot t)$,
and define $C_b$ likewise. Expanding the left side and using the bijection
$(x,h)\mapsto(x,y=x+h)$ gives
\[
 \E_t(\E_y C_a(y;t))(\E_x C_b(x;t)).
\]
Its absolute value is at most
$\bigl(\E_t|\E_yC_a(y;t)|^2\bigr)^{1/2}
 \bigl(\E_t|\E_xC_b(x;t)|^2\bigr)^{1/2}$.
Each squared expression is the $(d+1)$-cube average: expand the square
and substitute its second base point as the first plus a new increment.
It therefore equals the corresponding $U^{d+1}$ norm to the power
$2^{d+1}$. This proves the inequality. The original left side is
nonnegative because it is an average of powers of Gowers seminorms.
\end{proof}

Choose a cyclic group $Q$ with size between $C_sN$ and $2C_sN$ for a fixed
sufficiently large $C_s$, and zero-extend all functions supported on $[N]$.
The normalized interval and group norms differ by the fixed indicator
factor. For each fixed order that factor is bounded above by one and below
by $N/|Q|$, using the finite-group Gowers--Cauchy--Schwarz inequality.
Thus a fixed positive interval norm gives a fixed positive group norm,
and conversely; the constants depend only on the fixed ambient ratio and
order. Taking $C_s$ larger than $2^{s+1}$ ensures the same ambient group
is valid for all interval orders used here. The zero-extended product in
\eqref{eq:factorcorr} is exactly the group product for $|h|\leq N$;
there is no wrapped overlap since the ambient size is larger than $4N$.

By interface 6, \eqref{eq:factorcorr} has a uniformly positive
$U^{s-1}[N]$ norm for a dense set of $h$. Let $a=f\theta$ and
$b=\overline f\theta'$, both zero-extended. The average over all group
shifts of $\|a(\cdot+h)b\|_{U^{s-1}(Q)}^{2^{s-1}}$ is bounded below by a
fixed positive constant: the contributing set of interval shifts has
positive density also in $Q$, and the other summands are nonnegative.
Lemma~\ref{lem:CS} and the fixed bound on $b$ imply
$\|a\|_{U^s(Q)}$ is positive by a standard amount, hence so is
$\|f\theta\|_{U^s[N]}$. The inverse theorem at order $s-1$ now makes
$f\theta$ correlate with a degree-at-most-$s-1$ nilsequence $\psi$.
The actual product $\theta\psi$ is a degree-at-most-$s$ nilsequence on
the product filtered nilmanifold. Thus $f$ is $s$-biased.

\section{Starting and completing the induction}

The bounded induction base is the separately frozen and reviewed
\texttt{gi-base-002}. It starts from the exact forward bracket theorem,
proves the affected trivial-subgroup specialization, constructs a
periodic connected polynomial test and transfers its correlation
directly to the normalized interval. The older base in assembly001
is superseded and remains historical. No unpublished correction is
treated as an accepted black box.

For completeness the initial derivative step has the same normalization
control. The cube identity on the fixed ambient cyclic group gives
\[
 \E_h\|\Delta_h f\|_{U^s(Q)}^{2^s}
 =\|f\|_{U^{s+1}(Q)}^{2^{s+1}}.
\]
If the interval norm is positive by a standard amount, so is the right
side. Since $f$ is bounded, all summands have a fixed upper bound. The
nonzero differences occur only on the shifts represented by $[-N,N]$,
so a dense subset of those shifts has a uniformly positive $U^s[N]$
norm. Apply the inverse theorem of order $s-1$, then interface 2, to
obtain the family required by interface 3. Interfaces 3--5 and the preceding
section complete the induction step. The bounded limit-$N$ case is included:
a bounded interval has standard finite size, and any function on it is
represented by a Lipschitz function along a rational circle rotation,
with a standard finite Lipschitz bound. Thus the induction proves the
limit polynomial inverse conclusion for every limit $N$, conditional on
the listed intermediate interfaces and base.

We spell out the finite-family extraction rather than replace it by
pointwise compactness. Enumerate the rational filtered nilmanifolds of
degree at most $s$ up to isomorphism as $M_1,M_2,\ldots$, with one fixed
smooth metric on each, using interface 7. If the finite polynomial inverse
statement failed at some fixed $\delta>0$, for each $j$ choose a
$1$-bounded $f_j:[N_j]\to\C$ with norm at least $\delta$, but correlation
at most $1/j$ against every polynomial nilsequence from the first $j$
manifolds with $1$-bounded test and Lipschitz norm at most $j$.
Its ultralimit has positive norm, so the limit inverse conclusion gives a
fixed standard filtered nilmanifold, bounded Lipschitz limit test and
positive standard correlation. Normalize the test by a fixed bound to
make it $1$-bounded. It is isomorphic to some fixed $M_{j_0}$; fixed
metric equivalence transfers its Lipschitz bound to a standard constant.
Representatives of the polynomial and test satisfy the same bounds with
one slightly larger standard constant on an ultrafilter-large set.
For $j$ larger than that constant and $j_0$, the selected correlation
upper bound tends to zero, contradicting the positive lower bound.
This yields the finite-family polynomial statement with all parameters
fixed independently of $N$.

Finally interface 8 converts every member of the finite polynomial family
to a fixed nilmanifold of step at most $s$, with a fixed Lipschitz bound
and exactly equal sequence. A finite maximum of those bounds and the
unchanged positive correlation give the exact linear target in section 1.
This final join depends on the reviewed increment construction; it does
not import the failed full-polynomial suspension claim.

\section{What remains before this is a complete proof}

The local units and their exact incoming connections are now assembled
relative to their retained published inputs. The affine coefficient transfer
and multivariable factorization, bias and torsion conclusions have separate
independent decisions. Their substitutions and the degree-zero converse
boundary are new explicit connections in this successor. The earlier
composition decision is not silently extended to these new bytes.

The retained published forward bracket/Bohr, quantitative equal-side
equidistribution, rational-coordinate and compact cube-constraint inputs
keep their exact source scope. Related-use coverage and the complete
affected-dependency audit remain separate. This packet is submitted for
independent acceptance of the assembled target and its actual current
consumer interfaces. It makes no release-incorporation or whole-source
approval claim. The target keeps its original normalized norm, quantifier
order, fixed finite family and same-step linear output.

\end{document}
